Three fracture modes were assigned automatically (Sec.~\ref{sec:damage}): single-avalanche brittle failure (26 of 85 Stage~1--6 designs), stepwise failure with a few discrete events (15) and progressive failure with many events and load retention (43). The entropy-based localisation $L$ separates \emph{crack-like} from \emph{distributed} damage but not abrupt from progressive: precracked and holed sheets have $L=0.44$--0.54 (the crack path occupies a few grid cells), whereas periodic meshes that fail in one avalanche spread that avalanche over all equivalent ligaments ($L\approx0.2$--0.35), so their damage is delocalised in space but not in strain (Fig.~\ref{fig:loc}). The temporal signature -- number of discrete damage steps and post-peak strain range -- is therefore the operative distinction: abrupt designs have $3.5$ damage steps on average and a post-peak strain range below one increment, progressive designs $11$ steps and 5--15\% of additional strain. Across the design space the mode is set by the load-transfer capacity after the first local failure (Stage~4): many parallel ligaments of moderate strength dispersion (fine meshes, jittered meshes, load-parallel-vein meshes, distributed vacancies) fail sequentially; few wide strips, clustered defects and single flaws fail in one event. Random networks combine a localised crack path with several steps (they fail thinnest-ligament first, then re-localise).
